% ============================================================
%  InitialProposalDev — Notebook Build Audit & Execution System
%  Software Development Proposal — LaTeX Version
%  Compile with: latexmk -pdf InitialProposalDev.tex
% ============================================================
\documentclass[11pt, a4paper]{article}

% ── Packages ────────────────────────────────────────────────
\usepackage[T1]{fontenc}
\usepackage[utf8]{inputenc}
\usepackage[margin=2.4cm, top=2.8cm, bottom=2.8cm]{geometry}
\usepackage{microtype}
\usepackage{parskip}
\usepackage{amsmath}
\usepackage{booktabs}
\usepackage{array}
\usepackage{tabularx}
\usepackage{hyperref}
\usepackage{xcolor}
\usepackage{enumitem}
\usepackage{float}
\usepackage{caption}

\setlist[description]{font=\normalfont\itshape}

% ── Colour definitions ─────────────────────────────────────
\definecolor{clrBlue}{HTML}{2563EB}
\definecolor{clrGray}{HTML}{6B7280}

% ── Column types ───────────────────────────────────────────
\newcolumntype{Y}{>{\raggedright\arraybackslash}X}
\newcolumntype{R}{>{\raggedleft\arraybackslash}p{3.5cm}}
\newcolumntype{L}{>{\raggedright\arraybackslash}p{3cm}}

% ── Metadata ────────────────────────────────────────────────
\hypersetup{
  colorlinks=true,
  linkcolor=clrBlue,
  urlcolor=clrBlue,
  citecolor=clrBlue
}

% ============================================================
\begin{document}

% ── Title Page ──────────────────────────────────────────────
\begin{center}
  \vspace*{2cm}
  {\Huge\bfseries Software Development Proposal}\\[0.5cm]
  {\LARGE InitialProposalDev}\\[0.3cm]
  {\large Notebook Build Audit \& Execution System}\\[2.5cm]

  \begin{tabular}{rl}
    \textbf{Prepared for:} & Research \& Data Science Teams \\
    \textbf{Prepared by:}  & [Author Name / Consulting Team] \\
    \textbf{Date:}         & July 7, 2026 \\
  \end{tabular}

  \vfill
  {\footnotesize\color{clrGray} Confidential — Commercial in Confidence}
  \vspace*{1cm}
\end{center}

\newpage

% ── 1. Executive Summary & Problem Statement ───────────────
\section{Executive Summary \& Problem Statement}

\subsection*{The Challenge}

Computational notebooks (Jupyter, \texttt{.ipynb}) have become the dominant medium
for data science and machine learning workflows. However, they introduce systemic
risks that compromise reproducibility, correctness, and auditability at every stage
of the lifecycle:

\begin{itemize}
  \item \textbf{Untrusted notebook execution} — notebooks often originate from external
    collaborators, open-source repositories, or AI-assisted code generation tools
    without any deterministic guarantee that they run safely or correctly.

  \item \textbf{Lack of deterministic builds} — ad hoc environment pinning, missing
    dependency declarations, and absent seed controls cause notebooks to produce
    different results across machines or sessions, undermining scientific validity.

  \item \textbf{No structured audit protocol} — teams lack a systematic, multi-pass
    methodology for reviewing notebooks before production deployment, publication,
    or merge. Existing code review practices are insufficient for the unique failure
    modes of notebook-based ML pipelines (data leakage through preprocessing ordering,
    cross-validation contamination, inseparable training/inference logic).
\end{itemize}

The \textbf{Notebook Build Audit \& Execution System} addresses these gaps by
implementing the Construction Framework and Audit Framework defined in the
\texttt{NotebookBuildAudit} specification --- a formal six-pass review protocol
and a three-phase construction discipline --- within a containerized, AI-augmented
platform.

\subsection*{The Objective}

Build a web-based platform that enables research and data science teams to:

\begin{enumerate}
  \item \textbf{Execute notebooks deterministically} inside Docker-sandboxed
    environments with pinned dependencies and reproducibility controls.

  \item \textbf{Audit notebooks systematically} using the six-pass review protocol
    (Structural Overview, Reproducibility Check, Data Integrity Review, ML
    Correctness Audit, Code Quality Review, Deployment Readiness) powered by an
    LLM orchestration layer supporting both local (Ollama/\texttt{llama.cpp}) and
    cloud (OpenAI) providers. Additional cloud providers (Anthropic, OpenCode)
    are deferred to a subsequent release.

  \item \textbf{Interact with notebooks programmatically} via \texttt{@jupyter-kit}
    for cell-level inspection, execution, and modification.

  \item \textbf{Receive real-time audit logs and risk scores} (Low / Moderate /
    High per pass) through a modern frontend dashboard, closing the feedback loop
    between the Construction and Audit frameworks.
\end{enumerate}

% ── 2. Project Scope & Requirements ─────────────────────────
\section{Project Scope \& Requirements}

Based on the theoretical foundations in \texttt{NotebookBuildAudit.md} and
\texttt{NotebookBuildAudit.tex}, the software must include the following core
features:

\subsection*{Notebook Construction Workbench}

\begin{itemize}
  \item Section scaffold generator that pre-populates notebooks with the canonical
    eight-section header structure (Environment \& Dependencies, Configuration \&
    Global Parameters, Data Ingestion, Preprocessing \& Feature Engineering, Model
    Definition \& Training, Evaluation \& Metrics, Artifact Export, Conclusions \&
    Next Steps).

  \item Inline reproducibility controls: random seed configuration, deterministic
    flag toggles, and device configuration panels.

  \item Versioned artifact export routing: all persisted outputs (models, plots,
    reports) are routed through a single, versioned export convention.
\end{itemize}

\subsection*{Six-Pass Audit Engine}

\begin{description}
  \item[Pass 1 — Structural Overview] Produces a section map, red flag list, and
    recorded focus-area scope.

  \item[Pass 2 — Reproducibility Check] Generates a risk score (Low / Moderate /
    High) based on dependency pinning, seed configuration, hardcoded paths, and
    end-to-end re-executability.

  \item[Pass 3 — Data Integrity Review] Delivers a pipeline integrity report covering
    split ordering, pipeline-level leakage, missing-data handling, and
    ingestion-time validation.

  \item[Pass 4 — ML Correctness Audit] Produces a correctness checklist for
    evaluation metrics, cross-validation integrity, hyperparameter tuning boundaries,
    and baseline comparison.

  \item[Pass 5 — Code Quality Review] Outputs a code smell report flagging
    repetition ($\ge 3$-block threshold), dead code, naming quality, and output
    hygiene.

  \item[Pass 6 — Deployment Readiness] Returns a risk score (Low / Moderate / High)
    evaluating artifact export, inference/training separability, resource
    documentation, environment completeness, and data privacy.
\end{description}

\subsection*{LLM Orchestration Layer (Three-Level Audit)}

\begin{description}
  \item[Level 1 — Conceptual (Passes 1--2)] Assesses whether the notebook is
    structurally coherent and re-runnable.

  \item[Level 2 — Methodological (Passes 3--4)] Evaluates scientific soundness of
    the data pipeline and ML methodology.

  \item[Level 3 — Implementation (Passes 5--6)] Assesses code maintainability and
    production suitability.
\end{description}

\noindent Hybrid provider strategy: local LLMs via Ollama API or \texttt{llama.cpp}
for offline/privacy-sensitive evaluations; cloud providers (OpenAI)
for high-capability inference, routed through a provider-agnostic strategy
pattern that selects the appropriate model based on audit pass, latency requirements,
and privacy constraints. Additional cloud providers (Anthropic, OpenCode) are
supported by the extensible strategy pattern but deferred to a subsequent release.

\subsection*{@jupyter-kit Integration}

\begin{itemize}
  \item Programmatic notebook interaction: kernel management, cell-level execution,
    output capture, and metadata inspection.
  \item Sandbox execution pipeline: each notebook runs inside an isolated Docker
    container with pinned dependencies, controlled by \texttt{@jupyter-kit} cell
    execution commands.
\end{itemize}

\subsection*{Real-Time Dashboard}

\begin{itemize}
  \item Audit session management: submit notebooks, select pass scope (full six-pass
    or user-specified subsets), and configure LLM provider routing.
  \item Live audit progress: streaming per-pass deliverables as the LLM processes
    each level.
  \item Risk score visualization: Low / Moderate / High indicators per pass with
    drill-down to specific flagged cells.
\end{itemize}

\subsection*{Out of Scope}

\begin{itemize}
  \item Integration with proprietary CI/CD platforms beyond Docker-based deployment.
  \item Support for non-Python kernel languages (R, Julia) in the initial release.
  \item Automated code refactoring or rewrite suggestions --- the system is
    diagnostic, not prescriptive.
  \item Multi-tenant SaaS hosting; initial deployment targets single-instance,
    per-team Docker Compose environments.
  \item Custom model fine-tuning for the LLM audit layer.
\end{itemize}

% ── 3. Proposed Technology Stack ────────────────────────────
\section{Proposed Technology Stack}

To ensure scalability, security, and high performance, I propose the following
technology stack for this project:

\begin{table}[H]
  \centering
  \caption{Technology Stack Summary}
  \label{tab:techstack}
  \renewcommand{\arraystretch}{1.4}
  \begin{tabularx}{\textwidth}{l Y}
    \toprule
    \textbf{Layer} & \textbf{Technology \& Rationale} \\
    \midrule
    Frontend &
    Vite + React 18 + TypeScript. Fast HMR (Hot Module Replacement), type-safe
    component development, optimal bundle splitting for real-time dashboard updates.
    State management via React Context or Zustand for audit session state. Styling
    via Tailwind CSS or CSS Modules. \\

    Backend &
    Python 3.11+ + FastAPI. Async-first REST API for handling streaming LLM audit
    responses (Server-Sent Events), fast request validation via Pydantic models,
    auto-generated OpenAPI documentation. \\

    AI / LLM Orchestration &
    Ollama API, \texttt{llama.cpp}, OpenAI SDK.
    Provider-agnostic strategy pattern with unified \texttt{LLMProvider} interface
    (extensible for additional providers post-launch).
    Supports local (Ollama/llama.cpp) and cloud (OpenAI)
    routing. Prompt templates from \texttt{NotebookBuildAudit.tex} (Levels 1--3)
    are parameterized and injected per audit pass. \\

    Notebook Execution &
    \texttt{@jupyter-kit}. Programmatic Jupyter kernel management, cell execution,
    output capture, and notebook metadata manipulation within the sandboxed Docker
    environment. \\

    Database &
    PostgreSQL 15 + SQLAlchemy 2.0 (async). Relational storage for audit sessions,
    notebook metadata, audit results, and user configuration. Migration management
    via Alembic. \\

    Infrastructure &
    Docker + Docker Compose. Fully containerized, deterministic execution
    environments. Each notebook audit runs in an isolated Docker container with a
    pinned Python environment. Docker Compose orchestrates the full stack (FastAPI,
    PostgreSQL, React frontend, and per-audit sandbox containers). \\

    \bottomrule
  \end{tabularx}
\end{table}

% ── 4. High-Level Architecture ──────────────────────────────
\section{High-Level Architecture}

\subsection*{Conceptual Diagram}

\begin{verbatim}
[React + TypeScript Dashboard]  <-->  [FastAPI REST + SSE Gateway]
                                          |
                    +---------------------+---------------------+
                    |                     |                     |
              LLM Orchestrator    Audit State Machine    Jupyter Execution Engine
              (Provider Strategy)  (Six-Pass Protocol)   (@jupyter-kit Sandbox)
                    |                     |                     |
              +-----+-----+           Audit Results DB        Docker Sandbox
              |           |           (PostgreSQL)            (per-notebook container)
           Ollama      OpenAI
           (local)     (cloud)
\end{verbatim}

A detailed technical architecture diagram, API route specifications, and data flow
diagrams will be provided upon project commencement as part of the initial
deliverables.

\subsection*{Key Architectural Decisions}

\begin{description}
  \item[Strategy Pattern for LLM Providers] A \texttt{LLMProvider} abstract base
    class defines \texttt{generate(prompt, **kwargs)}. Concrete implementations
    (\texttt{OllamaProvider}, \texttt{OpenAIProvider})
    handle provider-specific API formatting. The
    \texttt{LLMOrchestrator} service selects the appropriate provider at runtime
    based on the audit configuration. The strategy pattern is designed to accept
    additional providers (Anthropic, OpenCode) post-launch without architectural
    changes.

  \item[Audit State Machine] Each notebook audit is modeled as a state machine
    transitioning through \texttt{Scope -> Pass1 -> Pass2 -> ... -> Pass6}. Gate
    decisions between Levels 1, 2, and 3 are rendered as configurable checkpoints
    --- the system can halt at a blocking issue or continue to surface further
    findings.

  \item[Docker Sandbox per Notebook] For each audit session, the system spawns a
    short-lived Docker container with the notebook's declared dependencies
    pre-installed. The container runs a headless Jupyter kernel managed by
    \texttt{@jupyter-kit}. The container is destroyed after the audit completes,
    providing full isolation.

  \item[Server-Sent Events (SSE) for Real-Time Audit Logs] The FastAPI backend
    streams per-pass audit deliverables to the React frontend via SSE, enabling
    live progress visualization without polling.
\end{description}

% ── 5. Project Deliverables ─────────────────────────────────
\section{Project Deliverables}

Upon completion, the client will receive:

\begin{itemize}
  \item Fully functional web application (Docker Compose deployable).
  \item Complete source code repository (transferred via GitHub).
  \item Database schema and Alembic migration scripts.
  \item REST API documentation (auto-generated OpenAPI/Swagger).
  \item Basic user documentation covering audit session creation, LLM provider
    configuration, and pass scope selection.
  \item Docker setup guide with per-notebook sandbox configuration.
\end{itemize}

% ── 6. Timeline & Milestones ────────────────────────────────
\section{Timeline \& Milestones}

The estimated time to complete this project is \textbf{8 weeks (2 months)}, broken down into
the following phases:

\begin{table}[H]
  \centering
  \caption{Project Phases}
  \label{tab:phases}
  \renewcommand{\arraystretch}{1.4}
  \begin{tabularx}{\textwidth}{c Y c}
    \toprule
    \textbf{Phase} & \textbf{Description} & \textbf{Duration} \\
    \midrule
    Phase 1 &
    \textbf{Environment \& Dockerization.} Initialize monorepo structure (backend,
    frontend, shared schemas). Create Docker Compose configuration with FastAPI,
    PostgreSQL, and frontend services. Implement per-notebook Docker sandbox
    spawning with \texttt{@jupyter-kit} kernel management. Set up dependency
    pinning and environment export capture. &
    2 weeks \\

    Phase 2 &
    \textbf{FastAPI Backend \& LLM Integration.} Implement FastAPI REST API with
    Pydantic models for audit sessions, notebook upload, and pass configuration.
    Build \texttt{LLMProvider} strategy pattern with Ollama and OpenAI
    concrete implementations (extensible interface ready for additional providers).
    Develop the \texttt{AuditStateMachine}
    implementing the six-pass protocol with gate decisions between Levels 1--3.
    Integrate \texttt{@jupyter-kit} for notebook parsing, cell execution, and
    output capture. Implement SSE streaming for real-time audit progress. &
    3 weeks \\

    Phase 3 &
    \textbf{Frontend Interface.} Scaffold Vite + React + TypeScript project. Build
    audit session dashboard (notebook upload, provider routing configuration, pass
    scope selection). Implement live audit progress view with per-pass risk score
    visualization and flagged-cell drill-down. Develop
    LLM provider configuration panel (local Ollama endpoint, cloud OpenAI API key). &
    2 weeks \\

    Phase 4 &
    \textbf{Integration \& QA.} End-to-end integration testing with real notebooks
    across all six passes. Validate LLM prompt templates against reference outputs
    from \texttt{NotebookBuildAudit.tex}. Security review of Docker sandbox isolation
    and API authentication. Final documentation and deployment guide. &
    1 week \\

    \bottomrule
  \end{tabularx}
\end{table}

% ── 7. Pricing & Engagement Model ───────────────────────────
\section{Pricing \& Engagement Model}

\subsection*{Engagement Terms}

\begin{itemize}
  \item \textbf{Team:} 2 dedicated developers allocated exclusively to this project.
  \item \textbf{Developer Rate:} \$1{,}000\ \text{USD per developer per month}
    ($\approx 3{,}900{,}000$ COP).
  \item \textbf{Communication:} Weekly syncs every Monday via Google Meet.
\end{itemize}

\subsection*{Investment}

\begin{table}[H]
  \centering
  \caption{Total Cost Breakdown}
  \label{tab:cost}
  \renewcommand{\arraystretch}{1.3}
  \begin{tabular}{l r r}
    \toprule
    \textbf{Cost Item} & \textbf{USD} & \textbf{COP (approx.)} \\
    \midrule
    Developer Costs (2 devs $\times$ 2 months $\times$
    \$1{,}000/mo) & \$4{,}000 & $\approx$ 15{,}600{,}000 COP \\
    Hardware / Server PC & \$1{,}000 & $\approx$ 3{,}900{,}000 COP \\
    \midrule
    \textbf{Total Project Cost} & \textbf{\$5{,}000} & \textbf{$\approx$
    19{,}500{,}000 COP} \\
    \bottomrule
  \end{tabular}
\end{table}

\subsection*{Payment Schedule}

\begin{enumerate}
  \item \textbf{30\% Upfront deposit} to start development
    ($\approx$ \$1{,}500 USD / $\approx$ 5{,}850{,}000 COP).

  \item \textbf{40\% Upon completion of Phase 2} --- Backend \& LLM Integration
    ($\approx$ \$2{,}000 USD / $\approx$ 7{,}800{,}000 COP).

  \item \textbf{30\% Upon final delivery and deployment}
    ($\approx$ \$1{,}500 USD / $\approx$ 5{,}850{,}000 COP).
\end{enumerate}

% ── 8. Maintenance & Support ────────────────────────────────
\section{Maintenance \& Support}

Following the final deployment, this proposal includes \textbf{2 months} of
complementary technical support.

\begin{description}
  \item[Included:] Bug fixes, security patches, and minor UI adjustments.
  \item[Not included:] Development of new features (these can be quoted separately
    or handled via a monthly retainer after the complementary period ends).
\end{description}

% ── 9. Next Steps & Acceptance ──────────────────────────────
\section{Next Steps \& Acceptance}

To proceed with this proposal, please sign below or confirm via email. Upon
acceptance, I will send over the formal contract and the initial invoice.

\vspace{1.5cm}

\noindent
\textbf{Accepted by:} \makebox[5cm]{\dotfill} \hspace{1cm}
\textbf{Date:} \makebox[3cm]{\dotfill}

\end{document}
